\documentclass[11pt]{article}
\usepackage[margin=0.82in]{geometry}
\usepackage{amsmath,amssymb,amsthm}
\usepackage{graphicx}
\usepackage{microtype}
\usepackage{xcolor}
\usepackage{hyperref}
\hypersetup{colorlinks=true,linkcolor=blue!55!black,urlcolor=blue!55!black,citecolor=blue!55!black}
\newtheorem{theorem}{Theorem}
\newtheorem{proposition}[theorem]{Proposition}
\newtheorem{corollary}[theorem]{Corollary}
\newtheorem{lemma}[theorem]{Lemma}
\newtheorem{remark}[theorem]{Remark}
\newcommand{\R}{\mathbb R}
\newcommand{\Sph}{\mathbb S}
\newcommand{\B}{\mathbb B}
\newcommand{\G}{\mathbb G}
\newcommand{\Hh}{\mathcal H}
\newcommand{\Lip}{\operatorname{Lip}}
\newcommand{\len}{\operatorname{length}}
\newcommand{\dist}{\operatorname{dist}}
\title{A full affirmative answer to the Grushin-plane\\no-axis-segment filling question}
\author{Candidate full-result packet for arXiv:1109.4641}
\date{27 August 2026}

\begin{document}
\maketitle

\begin{center}
\fbox{\parbox{0.92\linewidth}{\textbf{Status:} candidate full result,
likely valid.  Every Lipschitz loop in the first Grushin plane whose image
encloses no nondegenerate segment of the singular axis has a Lipschitz
filling.  Thus the result gives an affirmative answer to Question~8.7 of the
source.}}
\end{center}

\begin{abstract}
We prove the remaining implication by unfolding the finite-length continuum
traced by the loop.  Whyburn's cyclic elements are its maximal cycle blocks.
The no-enclosed-segment hypothesis forces every such block into one closed
Grushin half-plane.  The closed half-planes are CAT(0).  Replacing each cycle
block by a fresh copy of its half-plane, while retaining the dendritic
connectors, produces a complete CAT(0) tree of spaces.  The boundary loop
lifts to that tree of spaces with controlled Lipschitz constant and the
unfolding projects $1$-Lipschitzly to the Grushin plane.  Kirszbraun extension
in the CAT(0) unfolding and projection complete the proof.
\end{abstract}

\section{The source question and the result}

The first Grushin plane $G$ has underlying set $\R^2$, singular axis
\[
 Y=\{(0,y):y\in\R\},
\]
and, away from $Y$, length element
\begin{equation}\label{eq:metric}
 ds^2=dx^2+x^{-2}dy^2.
\end{equation}
Equivalently, its orthonormal horizontal vector fields are $\partial_x$ and
$x\partial_y$.  We write $d_G$ for its Carnot--Caratheodory distance.

DeJarnette--Hajlasz--Lukyanenko--Tyson prove that every $L$-Lipschitz
loop avoiding $Y$ has an $L$-Lipschitz filling.  They exhibit a bi-Lipschitz
circle enclosing a nondegenerate segment of $Y$ which has no Lipschitz
filling, and ask~\cite[Question~8.7]{DHJT}:

\begin{center}
\includegraphics[width=0.97\textwidth]{figures/open_problem_crop.png}

\smallskip
\emph{Question 8.7 on PDF page 34 of the source.}
\end{center}

We use the natural precise meaning of ``encloses'': a set $K\subset\R^2$
encloses a segment if a nondegenerate segment of $Y$ is contained in a
bounded component of $\R^2\setminus K$.  This is the meaning used by the
source example.

\section{Idea of the proof}

The obstruction is not the number of times a loop meets the singular axis.
It is whether a genuine cycle of its image uses both sides of that axis.  The
image is a finite-length Peano continuum, so Whyburn's cyclic-element theory
splits it into maximal cycle blocks joined by dendritic connectors.  A block
that met both open half-planes would contain a cross-axis Jordan curve; the
Jordan domain then forces the full image to enclose an axis interval.
Therefore every block lies in one closed Grushin half-plane.

Each closed half-plane is CAT(0).  We unfold the image by giving every cycle
block its own copy of the appropriate half-plane and gluing the copies to the
dendritic connectors only at their cut points.  This is a tree gluing along
convex singletons, so its completion is CAT(0).  Boundary subarcs lift without
length increase, and the unfolding projects $1$-Lipschitzly back to $G$.
Kirszbraun extension in the unfolded CAT(0) target followed by this projection
gives the filling.

\begin{theorem}[full answer to Question 8.7]\label{thm:main}
Let $\varphi:\Sph^1\to G$ be Lipschitz and put
$K=\varphi(\Sph^1)$.  If no nondegenerate segment of $Y$ lies in a bounded
component of $\R^2\setminus K$, then $\varphi$ extends to a Lipschitz map
$F:\B^2\to G$.

If the unit circle is given its Euclidean chord metric and
$\Lip(\varphi)=L$, one may take $\Lip(F)\le (\pi/2)L$.
\end{theorem}

The proof has three steps.  First, each closed Grushin half-plane is CAT(0).
Second, every genuine cycle block of $K$ lies in one of those half-planes.
Third, replacing cycle blocks by copies of their half-planes gives a CAT(0)
unfolding to which the loop lifts.

\section{Closed Grushin half-planes are CAT(0)}

Set
\[
 G_+=\{(x,y)\in G:x\ge0\},\qquad
 G_-=\{(x,y)\in G:x\le0\}.
\]
They carry the distance restricted from $G$.  Folding $x$ to $|x|$ preserves
the length of every rectifiable Grushin curve whose endpoints are in $G_+$.
Consequently the restricted distance on $G_+$ equals its intrinsic length
distance.  The corresponding statement holds for $G_-$.

\begin{lemma}[half-plane CAT(0) lemma]\label{lem:cat-half}
The complete geodesic spaces $G_+$ and $G_-$ are CAT(0).
\end{lemma}

\begin{proof}
It is enough to treat $G_+$.  For $\varepsilon>0$, translate the smooth
region
\[
 N_\varepsilon=\{(x,y):x\ge\varepsilon\}
\]
back to the fixed set $G_+$.  Its pulled-back Riemannian metric is
\begin{equation}\label{eq:epsmetric}
 g_\varepsilon=dx^2+(x+\varepsilon)^{-2}dy^2,
\end{equation}
and we denote its distance by $d_\varepsilon$.

The Gaussian curvature in the positive smooth Grushin region is
$-2x^{-2}\le0$.  The region $N_\varepsilon$ is complete and locally convex.
For the latter assertion, the $x$-coordinate of a geodesic satisfies
\[
 x''=-x^{-3}(y')^2\le0;
\]
thus a geodesic segment with endpoints in $x\ge\varepsilon$ remains there.
The metric Cartan--Hadamard theorem, or equivalently the standard convex
boundary form of it, shows that $(G_+,d_\varepsilon)$ is CAT(0).

As $\varepsilon\downarrow0$, the metrics in \eqref{eq:epsmetric} increase,
so $d_\varepsilon$ increases pointwise and is bounded above by $d_G$.
In fact
\begin{equation}\label{eq:metriclimit}
 d_\varepsilon(p,q)\uparrow d_G(p,q),\qquad p,q\in G_+.
\end{equation}
Here is a direct compactness check.  Take constant-speed almost minimizers
$\gamma_\varepsilon=(x_\varepsilon,y_\varepsilon)$ for
$d_\varepsilon(p,q)$.  A uniform length bound controls $x_\varepsilon$ and,
on the resulting bounded $x$-range, also controls $y_\varepsilon$.  After
passing to a uniformly convergent subsequence, write
\[
 b_\varepsilon=\frac{y_\varepsilon'}{x_\varepsilon+\varepsilon}.
\]
The controls $(x_\varepsilon',b_\varepsilon)$ are bounded in the appropriate
$L^2$ space.  Weak compactness gives $y'=xb$ for the limit curve, including
$y'=0$ almost everywhere on $\{x=0\}$, and lower semicontinuity gives
\[
 \len_G(\gamma)\le
 \liminf_{\varepsilon\downarrow0}\len_{g_\varepsilon}
 (\gamma_\varepsilon).
\]
This proves \eqref{eq:metriclimit}.

Pass now to the CAT(0) midpoint inequality.  For fixed $p,q,z$, let
$m_\varepsilon$ be the $d_\varepsilon$-midpoint of $p,q$.  The preceding
compactness argument gives a subsequence converging to a $d_G$-midpoint $m$.
On every relevant Euclidean compact set, \eqref{eq:metriclimit} is uniform:
the distance functions increase continuously to the continuous limit, so
Dini's theorem applies.  Hence
\[
 d_G(z,m)^2\le
 \frac12d_G(z,p)^2+\frac12d_G(z,q)^2-\frac14d_G(p,q)^2.
\]
The space $G_+$ is complete (indeed it is closed in the proper Grushin
plane), and the midpoint inequality characterizes complete CAT(0) spaces.
Reflection proves the claim for $G_-$.
\end{proof}

Two immediate consequences will be used below.  First, the identity map from
either $G_+$ or $G_-$ into $G$ is an isometric embedding.  Second, the
Lang--Schroeder Kirszbraun theorem applies to maps from subsets of Euclidean
space into either half-plane~\cite{LangSchroeder}.

\section{Finite-length continua and their cycle blocks}

We record the classical continuum-theoretic input and its quantitative
metric consequence.  A \emph{true cyclic element} of a Peano continuum $X$
is a maximal nondegenerate subcontinuum $C$ such that every two points of $C$
lie on a Jordan curve contained in $C$.

We use the following standard facts from cyclic element theory
\cite[Chapter~IV]{Whyburn}.
\begin{enumerate}
\item Every Jordan curve in $X$ is contained in a true cyclic element.
\item Distinct true cyclic elements meet in at most one point, necessarily a
cut point of $X$.
\item If $C$ is a true cyclic element, the closure of each component of
$X\setminus C$ meets $C$ in at most one point.
\item After the true cyclic elements are regarded as blocks, the remaining
connectors are dendritic and the block-incidence structure is a tree: there
is a unique simple cyclic chain between any two elements.
\end{enumerate}
For orientation, the point-attachment assertion in (3) is forced by
maximality.  If a component outside $C$ attached at distinct points $a,b$,
local connectedness would supply an outside arc from $a$ to $b$.  Combining
it with a suitable arc in a Jordan curve of $C$ would enlarge the cyclic
element.

The next lemma packages exactly the infinite gluing statement needed here.

\begin{lemma}[CAT(0) cyclic unfolding]\label{lem:unfold}
Let $M$ be a complete metric space and let $X\subset M$ be a compact
connected set with $\Hh^1_M(X)<\infty$.  For every true cyclic element
$C\subset X$, suppose there are
\begin{itemize}
\item a complete CAT(0) space $Z_C$,
\item a map $i_C:C\to Z_C$, and
\item a $1$-Lipschitz map $p_C:Z_C\to M$
\end{itemize}
such that $p_Ci_C$ is the inclusion $C\hookrightarrow M$, and
\begin{equation}\label{eq:blocklength}
 \len_{Z_C}(i_C\circ\alpha)\le\len_M(\alpha)
\end{equation}
for every rectifiable path $\alpha$ in $C$.

Then there are a complete CAT(0) space $E$, a $1$-Lipschitz map
$P:E\to M$, and a lift $j:X\to E$ with $Pj$ equal to the inclusion, such that
for every rectifiable path $\alpha$ in $X$,
\begin{equation}\label{eq:globallength}
 \len_E(j\circ\alpha)\le\len_M(\alpha).
\end{equation}
\end{lemma}

\begin{proof}
A compact connected set of finite $\Hh^1$-measure is locally connected and
rectifiably connected, hence is a Peano continuum.  Apply its cyclic-element
decomposition.  Replace every true cyclic element $C$ by a disjoint copy of
$Z_C$.  Retain each dendritic connector with its intrinsic arc-length metric,
which is an $\R$-tree metric, and identify every attachment point with the
corresponding point $i_C(q)$ in the adjacent block.  Distinct blocks and
connectors are therefore glued only along singletons.  The tree property in
(4) says that no cycle is introduced between the pieces.

For completeness, root the cyclic-element tree.  Its nondegenerate block
traces and nondegenerate connector components are countable: their interiors
are disjoint, each has positive $\Hh^1$-measure, and their total measure is at
most $\Hh^1(X)$.  The remaining degenerate elements are limit points of this
system.  Take an increasing sequence of finite rooted subtrees that exhausts
the nondegenerate pieces.  The point-attachment property ensures that every
piece newly exposed by a rooted subtree meets the preceding subtree in
exactly its parent attachment point.  If infinitely many pieces occur along
one rooted chain, first include those whose \emph{original traces in $X$}
have diameter at least $1/n$, then let $n\to\infty$; there are only finitely
many at each scale because $\Hh^1(X)<\infty$.  This gives an increasing,
dense union
\[
 E_1\subset E_2\subset\cdots
\]
in which each new piece is attached to the old union at one point.  Each
finite stage is CAT(0) by Reshetnyak gluing along the closed convex set
$\{q\}$, and its inclusion in the next stage is isometric and convex
\cite[Chapter~II.11]{BridsonHaefliger}.  Every geodesic triangle in the union
lies in some finite stage and satisfies the CAT(0) comparison inequality.
Let $E$ be its metric completion.  Completions of CAT(0) spaces are CAT(0),
so $E$ is complete CAT(0).  This is also the standard metric realization of
the cyclic-element tree with the blocks $C$ replaced by $Z_C$.

On $Z_C$ define $P=p_C$.  On every dendritic connector use its original map
into $X\subset M$.  These maps agree at attachment points.  They are
$1$-Lipschitz on pieces, so the path-metric definition makes $P$ short on the
dense union; completeness of $M$ extends it to a $1$-Lipschitz map on $E$.

The canonical lift $j$ sends a point of $C$ to $i_C(C)$ and a point in the
dendritic part to its corresponding tree point.  Shared cut points were
identified, so this is well defined and $Pj$ is the inclusion.  A rectifiable
path in $X$ lifts piece by piece.  On block passages its length does not
increase by \eqref{eq:blocklength}; on dendritic passages the intrinsic
arc-length is the original path length.  First sum over finitely many
passages and then exhaust the countably many passages.  Monotone convergence
of length gives \eqref{eq:globallength}.  In particular, this estimate makes
$j\circ\alpha$ continuous even at accumulating attachment points.
\end{proof}

\begin{remark}\label{rem:why-finite}
The finite-length hypothesis is doing two jobs: it makes $X$ a Peano
continuum, and it makes all nontrivial pieces countable with summable sizes.
Thus the completion in Lemma~\ref{lem:unfold} is not hiding a local-finiteness
assumption.
\end{remark}

\section{The no-segment condition makes every block one-sided}

The metric induced on the singular axis has the snowflake scale
\begin{equation}\label{eq:snowflake}
 d_G((0,s),(0,t))\asymp |s-t|^{1/2}.
\end{equation}
Consequently every nondegenerate Euclidean interval in $Y$ has infinite
$1$-dimensional Hausdorff measure for $d_G$.

\begin{lemma}[one-sided cyclic elements]\label{lem:onesided}
Let $K\subset G$ be compact, connected, and of finite $\Hh^1_G$-measure.
Assume that no nondegenerate segment of $Y$ is contained in a bounded
component of $\R^2\setminus K$.  Then every true cyclic element of $K$ is
contained in $G_+$ or in $G_-$.
\end{lemma}

\begin{proof}
Suppose a true cyclic element $C$ contains $a$ with positive $x$-coordinate
and $b$ with negative $x$-coordinate.  By the defining property of $C$ there
is a Jordan curve $J\subset C$ through $a$ and $b$.  Let $\Omega$ be its
bounded Jordan domain.

The set $\Omega$ meets $Y$.  Otherwise connectedness would put $\Omega$ in
one component of $\R^2\setminus Y$, and then
$J=\partial\Omega$ would lie in the corresponding closed half-plane,
contrary to the choice of $a,b$.  Since $\Omega$ is open,
$\Omega\cap Y$ contains a nondegenerate interval $I$.

The finite-measure set $K$ cannot contain an axis interval, by
\eqref{eq:snowflake}.  Thus the compact set $K\cap I$ has empty relative
interior in $I$, and $I\setminus K$ contains a nondegenerate open
subinterval $I_0$.  Because $I_0$ is connected, it lies in one component
$U$ of $\R^2\setminus K$.  The Jordan curve $J\subset K$ prevents $U$ from
leaving $\Omega$, so $U$ is bounded.  Hence $K$ encloses $I_0$, a
contradiction.  Therefore $C$ cannot meet both open half-planes, proving the
claim.
\end{proof}

\section{Proof of the full theorem}

\begin{proof}[Proof of Theorem~\ref{thm:main}]
The image $K=\varphi(\Sph^1)$ is compact and connected, and
\[
 \Hh^1_G(K)\le \len_G(\varphi)<\infty.
\]
By Lemma~\ref{lem:onesided}, every true cyclic element $C$ of $K$ is
contained in one of $G_+$ and $G_-$.  Choose one such half-plane and call it
$Z_C$.  Let $i_C:C\to Z_C$ and $p_C:Z_C\to G$ both be the coordinate
identity.  Lemma~\ref{lem:cat-half} says that $Z_C$ is complete CAT(0).
Folding shows that its intrinsic metric agrees with the distance restricted
from $G$, so $p_C$ is $1$-Lipschitz and
\[
 \len_{Z_C}(i_C\circ\alpha)=\len_G(\alpha)
\]
for every rectifiable path $\alpha$ in $C$.

Apply Lemma~\ref{lem:unfold}.  We obtain a complete CAT(0) space $E$, a short
map $P:E\to G$, and a lift $j:K\to E$ preserving the length upper bound on
every rectifiable path.  Put
\[
 \widetilde\varphi=j\circ\varphi:\Sph^1\to E.
\]
For $s,t\in\Sph^1$, let $A_{s,t}$ be the shorter circular arc between them.
The length estimate in Lemma~\ref{lem:unfold} gives
\[
\begin{split}
 d_E(\widetilde\varphi(s),\widetilde\varphi(t))
 &\le \len_E(\widetilde\varphi|_{A_{s,t}})\\
 &\le \len_G(\varphi|_{A_{s,t}})\\
 &\le L\,\len_{\R^2}(A_{s,t})
 \le \frac\pi2 L|s-t|.
\end{split}
\]
Thus $\widetilde\varphi$ is $(\pi/2)L$-Lipschitz on the unit circle with its
Euclidean chord metric.

Euclidean space has curvature bounded below by $0$, and $E$ is complete
CAT(0).  The Lang--Schroeder Kirszbraun theorem therefore extends
$\widetilde\varphi$ to a $(\pi/2)L$-Lipschitz map
\[
 \widetilde F:\B^2\to E.
\]
Finally set $F=P\circ\widetilde F$.  Since $P$ is $1$-Lipschitz,
$F$ has the same Lipschitz bound, and on the boundary
\[
 F|_{\Sph^1}=Pj\varphi=\varphi.
\]
This is the required filling.
\end{proof}

\section{Stress tests and novelty check}

The construction handles the configurations that defeat a direct
one-contact gluing argument.
\begin{itemize}
\item A figure eight with one lobe on each side of $Y$ becomes two CAT(0)
half-planes glued at the common cut point.
\item A theta block cannot mix the two sides: any two cross-side points in
the same true cyclic element lie on a cross-axis Jordan curve, contradicting
Lemma~\ref{lem:onesided}.
\item Infinitely many alternating one-sided blocks cause no completeness
problem.  Finite subtrees are CAT(0), and their metric completion is CAT(0).
\item A loop whose image has no Jordan curve at all lifts to an $\R$-tree,
regardless of how often it meets the singular axis.
\end{itemize}

A bounded novelty search through 27 August 2026 covered the run registry and
solution indexes; the exact wording of Question~8.7; searches combining
``Grushin plane'', ``Lipschitz extension'', ``cyclic element'', and
``CAT(0)''; and the later Grushin embedding and quasiconformal papers returned
by those searches.  No later paper claiming this answer or this unfolding
argument was found.  Novelty confidence is moderate rather than definitive,
because the proof combines classical ingredients in a short new way.

The packet is recommended for human review as a likely valid full answer.
The highest-value audit is Lemma~\ref{lem:unfold}: compare its tree
realization directly with Whyburn's cyclic-element decomposition, then verify
that the finite-stage exhaustion includes every accumulating connector.  The
finite-length argument in Remark~\ref{rem:why-finite} is included precisely
to expose that point.

\begin{thebibliography}{9}
\bibitem{DHJT}
N. DeJarnette, P. Hajlasz, A. Lukyanenko, and J. T. Tyson,
\emph{On the lack of density of Lipschitz mappings in Sobolev spaces with
Heisenberg target}, arXiv:1109.4641v3 (2014).
\url{https://arxiv.org/abs/1109.4641}.

\bibitem{LangSchroeder}
U. Lang and V. Schroeder,
\emph{Kirszbraun's theorem and metric spaces of bounded curvature},
Geom. Funct. Anal. 7 (1997), 535--560.

\bibitem{Whyburn}
G. T. Whyburn,
\emph{Analytic Topology}, American Mathematical Society Colloquium
Publications, vol. 28, 1942.

\bibitem{BridsonHaefliger}
M. R. Bridson and A. Haefliger,
\emph{Metric Spaces of Non-Positive Curvature}, Springer, 1999.
\end{thebibliography}

\end{document}
